\documentclass[11pt]{article}
\usepackage{amsmath}
\usepackage{amssymb}
\usepackage{geometry}
\geometry{margin=1in}

\title{MCT: Markov Chain Deriver/Generator/Tester \\ Project Specification}
\author{}
\date{}

\begin{document}
\maketitle

====================
\section{Core constraints / things to keep in mind}
====================

\begin{itemize}
\item Channels, not just Kraus operators  
\end{itemize}
  A \textit{valid} input is a completely positive trace-preserving (CPTP) map with Kraus form  
  
\begin{equation}
\rho' = \mathcal{E}(\rho) = \sum_k E_k \,\rho\, E_k^\dagger
\end{equation}

  and
  
\begin{equation}
\sum_k E_k^\dagger E_k = I.
\end{equation}

  Code should either enforce or check this normalization when given $E_k$.

\begin{itemize}
\item Basis is part of the specification  
\end{itemize}
  You must always carry an explicit ordered basis  
  
\begin{equation}
\mathcal{B} = \{\,|b_0\rangle, |b_1\rangle, \dots, |b_{d-1}\rangle \,\}
\end{equation}

  for the $d = 2^n$-dimensional Hilbert space. The Markov chain lives on these labels; changing basis changes the effective Markov dynamics.

\begin{itemize}
\item Dimensional blow-up and symbolic limits  
\end{itemize}
  For $n$ qubits, $\rho$ is $2^n \times 2^n$, and the superoperator is $4^n \times 4^n$.  
  Practical rule:  
\begin{itemize}
\item ``Fully symbolic, many-parameter'' mode: realistically only for $n \le 2$.  
\item Beyond that, numeric-only or \textit{very} restricted symbolic parameters.
\end{itemize}

\begin{itemize}
\item Distinguish exact quantum evolution vs population Markov chain  
\end{itemize}
  Exact evolution is the channel:
  
\begin{equation}
\rho_{t+1} = \mathcal{E}(\rho_t).
\end{equation}

  The \textit{classical} Markov chain is an \textit{approximate} description of populations in a chosen basis, defined by some rule such as:
\begin{itemize}
\item Start from basis projector $\rho_i = |b_i\rangle\langle b_i|$.
\item Evolve: $\rho'_i = \mathcal{E}(\rho_i)$.
\item Define transition probabilities
\end{itemize}
    
\begin{equation}
P_{ji} = \langle b_j | \rho'_i | b_j \rangle.
\end{equation}

  This gives a column-stochastic matrix $P$ when $\mathcal{E}$ is trace-preserving, but it discards coherence information by construction.

\begin{itemize}
\item Discrete step vs continuous-time generator  
\end{itemize}
  Decide which you are using as the \textit{primary} object:

\begin{itemize}
\item Discrete channel per time-step $\Delta t$:
\end{itemize}
    
\begin{equation}
\rho_{n+1} = \mathcal{E}_{\Delta t}(\rho_n).
\end{equation}

\begin{itemize}
\item Continuous-time Lindbladian generator $\mathcal{L}$:
\end{itemize}
    
\begin{equation}
\frac{d\rho}{dt} = \mathcal{L}(\rho), \quad
    \mathcal{E}_t = e^{t \mathcal{L}},
\end{equation}

    with $\mathcal{L}$ in Lindblad form
    
\begin{equation}
\mathcal{L}(\rho)
      = -\frac{i}{\hbar}[H,\rho]
        + \sum_j \Big(
          2 A_j \rho A_j^\dagger
          - A_j^\dagger A_j \rho
          - \rho A_j^\dagger A_j
        \Big).
\end{equation}

  Make the choice explicit in the API, because it affects how you interpret $T_1$, $T_2$, and per-step probabilities.

\begin{itemize}
\item Two backends from day one  
\item Symbolic backend (e.g.\ SymPy matrices) for small-$n$ analytic expressions.  
\item Numeric backend (NumPy / PyTorch) for realistic system sizes and for validating the Markov approximation against full density-matrix evolution.
\end{itemize}

====================
\section{High-level TODO list for the agents}
====================

Structure this as the module roadmap. Each bullet is a self-contained work package.

\begin{enumerate}
\item \textbf{Define core data structures}  
\begin{itemize}
\item A \texttt{QuantumChannel} object that stores:
\item Kraus operators $\{E_k\}$ (symbolic or numeric matrices) in a fixed basis.  
\item Dimension $d$ and number of qubits $n$.  
\item Optionally, the superoperator matrix $S$ acting on vectorized density matrices:
\end{itemize}
       
\begin{equation}
|\rho'\rangle\rangle = S \,|\rho\rangle\rangle.
\end{equation}

\begin{itemize}
\item A \texttt{Basis} object:
\item Explicit ordered list of computational basis states for $n$ qubits (\texttt{["00...0", ...]}).
\end{itemize}

\item \textbf{Implement ``apply channel to density matrix''}  
\begin{itemize}
\item Core function: \texttt{rho\_out = apply\_channel(channel, rho\_in)}  
\end{itemize}
       computing
       
\begin{equation}
\rho_{\text{out}} = \sum_k E_k \rho_{\text{in}} E_k^\dagger.
\end{equation}

\begin{itemize}
\item Should work in both:
\item symbolic mode (small systems, returns SymPy expressions), and  
\item numeric mode (returns NumPy or PyTorch arrays).
\end{itemize}

\item \textbf{Implement Markov-chain extraction from a channel}  
\begin{itemize}
\item Function: \texttt{P = make\_markov\_from\_channel(channel, basis)}  
\item Default rule (document this in the code!):
\item For each basis state $|b_i\rangle$:
\end{itemize}
       
\begin{equation}
\rho_i = |b_i\rangle\langle b_i|,\quad
       \rho'_i = \mathcal{E}(\rho_i).
\end{equation}

\begin{itemize}
\item Define
\end{itemize}
       
\begin{equation}
P_{ji} = \langle b_j | \rho'_i | b_j \rangle.
\end{equation}

\begin{itemize}
\item Verify in code:
\item $P_{ji} \ge 0$ and each column sums to $1$ (within tolerance).
\end{itemize}

\item \textbf{Implement a density-matrix simulator}  
\begin{itemize}
\item Function for \textit{discrete} channels: \texttt{rho\_t = simulate\_density\_discrete(rho0, channel, steps)}  
\end{itemize}
       repeated application:
       
\begin{equation}
\rho_{k+1} = \mathcal{E}(\rho_k).
\end{equation}

\begin{itemize}
\item Optional: function for \textit{continuous-time} generator $\mathcal{L}$:
\item either via exponentiation $\mathcal{E}_t = e^{t \mathcal{L}}$ for small systems, or numeric ODE integration of
\end{itemize}
       
\begin{equation}
\frac{d\rho}{dt} = \mathcal{L}(\rho).
\end{equation}


\item \textbf{Cross-check: Markov chain vs full density evolution}  
\begin{itemize}
\item Given $P$ and \texttt{simulate\_density\_discrete}, implement a test function:
\item Start from each basis projector $\rho_i = |b_i\rangle\langle b_i|$.  
\item Evolve with both:
\item full channel (density matrix), and  
\item Markov chain on populations.  
\item Compare
\end{itemize}
       
\begin{equation}
\text{diag}(\rho_k) \quad\text{vs}\quad P^k \, e_i,
\end{equation}

       where $e_i$ is the $i$-th standard basis vector.  
\begin{itemize}
\item Use this to quantify how ``classical'' the channel looks in a given basis.
\end{itemize}

\item \textbf{Build a library of standard test channels}  
   Implement constructors (both symbolic and numeric) for:

\begin{itemize}
\item Single-qubit amplitude damping (AD): Kraus operators
\end{itemize}
     
\begin{equation}
E_0 = \begin{pmatrix} 1 & 0 \\ 0 & \sqrt{1-\lambda} \end{pmatrix},\quad
     E_1 = \begin{pmatrix} 0 & \sqrt{\lambda} \\ 0 & 0 \end{pmatrix}.
\end{equation}

\begin{itemize}
\item Two-qubit amplitude damping via tensor products
\end{itemize}
     
\begin{equation}
K_0 = E_0\otimes E_0,\;
     K_1 = E_0\otimes E_1,\;
     K_2 = E_1\otimes E_0,\;
     K_3 = E_1\otimes E_1.
\end{equation}

\begin{itemize}
\item Single-qubit pure dephasing (phase damping) with Kraus operators like
\end{itemize}
     
\begin{equation}
K_0 = \sqrt{1-p}\,I,\quad K_1 = \sqrt{p}\,Z.
\end{equation}

\begin{itemize}
\item Simple depolarizing or Pauli channels.
\end{itemize}

   These give known steady states and known $T_1$, $T_2$ relations for validation.

\item \textbf{Validation and sanity checks}  
\begin{itemize}
\item CPTP check:
\item Verify
\end{itemize}
       
\begin{equation}
\sum_k E_k^\dagger E_k \approx I
\end{equation}

       numerically, or exactly in symbolic mode.  
\begin{itemize}
\item Markov matrix check:
\item $P_{ji} \ge 0$ and column sums $\sum_j P_{ji} = 1$.  
\item Physics sanity:
\item AD channel relaxes to ground state; dephasing leaves populations unchanged.  
\item For AD with rate $\gamma = 1/T_1$, verify population decay
\end{itemize}
       
\begin{equation}
\rho_{11}(t) \sim e^{-t/T_1}
\end{equation}

       and coherence decay consistent with chosen $T_2$.

\item \textbf{Document limitations explicitly}  
\begin{itemize}
\item State in the README/API:
\item Maximum recommended $n$ for symbolic mode.  
\item Default basis choice if none is supplied (e.g.\ computational).  
\item What Markov-construction rule is used and what information is lost (coherences).
\end{itemize}
\end{enumerate}

====================
\section{Minimal reference set for the agents}
====================

To keep them focused, refer to the following resources with indicated page ranges:

\begin{enumerate}
\item \textbf{\texttt{qubit\_guide.pdf} -- Chapter 9 (Quantum channels)}  
   Use this for the \textit{definitions and structure} of channels and Kraus operators.

\begin{itemize}
\item Pages 170--171: introduction to quantum channels, ``everything is secretly unitary'', and the idea that non-unitary evolution arises from unitary evolution + ancilla + partial trace.
\item Page 177: Stinespring dilation and derivation of the operator-sum form
\end{itemize}
     
\begin{equation}
\rho \mapsto \rho' = \sum_i E_i \rho E_i^\dagger
\end{equation}

     with $\sum_i E_i^\dagger E_i = I$.
\begin{itemize}
\item Page 178: Kraus representation explicitly stated and discussed; good for citing the core equation and completeness relation.
\item Page 181: composition of channels (adding ancilla, unitaries, partial trace) and the fact that each piece has its own Kraus representation; this motivates how you'll implement composition in code.
\end{itemize}

   That's enough from this book for the purposes of your project.

\item \textbf{\texttt{T1 T2 noise.pdf} -- amplitude damping example}  
   Use this as a \textit{worked example} of going from single-qubit AD Kraus operators to two-qubit Kraus operators and Bell-state behavior.

\begin{itemize}
\item Page 3: explicit AD Kraus operators
\end{itemize}
     
\begin{equation}
E_0 = \begin{pmatrix} 1 & 0 \\ 0 & \sqrt{1-\lambda} \end{pmatrix},\quad
     E_1 = \begin{pmatrix} 0 & \sqrt{\lambda} \\ 0 & 0 \end{pmatrix},
\end{equation}

     and construction of two-qubit Kraus operators by tensor products.  

   Agents can use this as a sanity-check case for their multi-qubit symbolic and numeric implementations.

\item \textbf{\texttt{Manenti\_Motta-Do\_not\_share.pdf} -- Markovian master equations, $T_1$, $T_2$}  
   Use this as the \textit{bridge} between Lindblad master equations, $T_1/T_2$, and the effective channels you are approximating.

\begin{itemize}
\item Page 300: general Markovian master equation in Lindblad form, with decay rates $\gamma_k$ and dissipators
\end{itemize}
     
\begin{equation}
\mathcal{D}[A_k](\rho) = A_k \rho A_k^\dagger - \frac{1}{2}\{A_k^\dagger A_k, \rho\},
\end{equation}

     and explicit amplitude-damping example linking $\gamma = 1/T_1$.
\begin{itemize}
\item Page 336: summary table including Markovian master equation, amplitude damping, and dephasing, plus the relations defining $T_1$ and $T_\phi$.
\end{itemize}

   This gives enough physics to relate your discrete-step channel parameters $(\lambda_1, \lambda_\phi)$ to continuous-time $T_1$, $T_2$.

\item \textbf{Optional / back-pocket: \texttt{Quantum Noise\_ ... Gardiner \& Zoller}}  
   Only if you want them to see the ``canonical'' Lindblad form in a more theoretical setting.

\begin{itemize}
\item Page 165: statement of the Lindblad form and discussion that it is the general generator of a quantum Markov process.
\end{itemize}

   You can mark this as ``for background only; not required for implementation.''
\end{enumerate}

If you want, the next step can be: turn this into a short, agent-facing spec document (with function names and signatures) that you can paste directly into an agentic coding session.

\end{document}